\documentclass[10pt,a4paper]{amsart}
\usepackage[margin=19mm]{geometry}
\usepackage{amsmath,amssymb,mathtools}
\usepackage[scr=rsfs,cal=euler]{mathalfa}
\usepackage{xfrac}
\usepackage[unicode,hidelinks,bookmarks=false]{hyperref}
\newcommand{\ie}{i.e.\ }
\newcommand{\cf}{cf.\ }
\newcommand{\resp}{respectively\ }
\newcommand{\Subsection}{\subsection}
\providecommand{\nobreakdash}{\nobreak\hbox{-}}
\setlength{\emergencystretch}{2em}
\multlinegap0pt
\usepackage{mathtools}
\usepackage{amssymb}
\usepackage{enumerate}
\usepackage[shortlabels]{enumitem}
\usepackage{float}
\usepackage{braket}
\usepackage{tikz}
\newtheorem{lemma}{Lemma}
\title{A central limit theorem for partitions involving generalised divisor functions}
\author{Madhuparna Das, Nicolas Robles}
\date{Source: arXiv:2510.19740v2 (CC BY 4.0)}
\begin{document}
\maketitle

\noindent\emph{Source and licence.} A central limit theorem for partitions involving generalised divisor functions. Version arXiv:2510.19740v2 (CC BY 4.0), distributed by arXiv under the Creative Commons Attribution 4.0 International licence, http://creativecommons.org/licenses/by/4.0/, as recorded in the arXiv licence metadata for this exact version and shown on the abstract page https://arxiv.org/abs/2510.19740v2. Full licence notice: \texttt{licenses/arxiv-2510.19740-CC-BY-4.0.txt}.\par\smallskip

\noindent\textit{Editorial scope.} Counted unit: the article's lemma shfiteddirichletseries (source0.tex lines 1041--1050). The article's complete original proof of it is delivered with it (source0.tex lines 1052--1076, one \texttt{proof} environment of 25 lines). No mathematics is written, repaired or re-derived: the statement and the proof are the article's own lines, in the article's own notation, and no retained line is substituted or re-flowed.  Retained uncounted as the source-local context the proof needs: lines 520--523; lines 539--542; lines 545--552; lines 585--611; lines 587--590; lines 781--808; lines 783--786.  Nothing was omitted from any retained span, so the package carries no omission record.  No retained line contains a citation, so the wrapper carries no bibliography and no bibliography entry.  No retained line contains a figure environment, an image-inclusion command or drawing code, so no image is delivered and the package declares no asset dependency.  Editorial formatting only, in addition: an AMS \texttt{amsart} preamble that restates the article's own declarations and macros used by the retained lines, namely the environments \texttt{lemma} on separate counters, together with the standard packages those lines need.  The retained environments print in this document as Lemma 1 in order of appearance, whereas the article prints the same items with the numbering of its own full document.  The complete native source of the article as distributed in the arXiv e-print for this exact version is bundled unchanged as \texttt{sources/187597/source0.tex}.
\begin{align}\label{kthdivisoridentity}
\sigma_r(n)
=\zeta(r+1)n^r\sum_{m=1}^{\infty}\frac{c_m(n)}{m^{r+1}},
\end{align} 

\begin{align}\label{generalizedgausssum}
c'_{\chi}(n) 
\coloneqq \sum_{\substack{b(\bmod{m})\\(b,m)=1}}\chi(b)e\left(\frac{bn}{m}\right).
\end{align}

\begin{lemma}\label{shiftedramanujansumformula}
Let $c_m(n)$ denote the Ramanujan sum for $m\geq 1$ and $n\geq 1$. Then the following identity holds
\begin{align*}
c_m(n+1)
=\frac{\mu(m)}{\varphi(m)}c_m(n)+\frac{1}{\varphi(m)}\sum_{\substack{\chi(\bmod{m})\\ \chi\ne\chi_0}}\tau(\chi)c'_{\bar{\chi}}(n),
\end{align*}
where $\tau(\chi)=c'_{\chi}(1)$ has defined in~\eqref{generalizedgausssum}. %Here, $\mu(m)$ denotes M\"{o}bius function, and $\varphi(m)$ is Euler’s totient function.
\end{lemma}

\begin{lemma}\label{dsklemma}
For a fixed integer $r\geq 1$, consider the double Dirichlet series
\begin{align}\label{dskdefinition}
D_1(s,r)
\coloneqq &\sum_{m=1}^{\infty}\frac{\mu(m)}{\varphi(m)m^{r+1}}\sum_{n=1}^{\infty}\frac{c_m(n)}{n^s},
\end{align}
where $c_m(n)$ denotes the Ramanujan sum. The following identity then holds for $\mathrm{Re}(s)>1$:
\begin{align*}
D_1(s,r)
=\frac{\zeta(s)}{\zeta(s+r+1)}K_r(s),
\end{align*}
where
\begin{align}\label{arsdefinition}
K_r(s)
=C(r)\prod_{p}\left(1+A_{p,r}(s)\right)
\end{align}
with 
\begin{align*}%\label{arpsdefinition}
A_{p,r}(s)
=-\frac{\frac{1}{p^{r+1}(p-1)}}{1+\frac{1}{p^{r+1}(p-1)}}\left(\frac{1}{p^s}-\frac{1}{p^{s+r+1}}+\frac{1}{p^{2s+r+1}}-\frac{1}{p^{2s+2r+2}}+\cdots-\frac{1}{p^{(k+1)(s+r+1)}}+\cdots\right)
\end{align*}
and
\begin{align}\label{crdefinition}
C(r)
\coloneqq\prod_{p}\left(1+\frac{1}{p^{r+1}(p-1)}\right).    
\end{align}
\end{lemma}    

\begin{align}
D_1(s,r)
\coloneqq &\sum_{m=1}^{\infty}\frac{\mu(m)}{\varphi(m)m^{r+1}}\sum_{n=1}^{\infty}\frac{c_m(n)}{n^s},
\end{align}

\begin{lemma}\label{gausssumlemma}
Let $r>1$ be fixed, and for a Dirichlet character $\chi$ modulo $m$, define
\begin{align}\label{d2sksum}
D_2(s,r)
\coloneqq\sum_{m=1}^{\infty}\frac{1}{\varphi(m)m^{r+1}}\sum_{\substack{\chi(\bmod{m})\\\chi\ne\chi_0}}\tau(\chi)\sum_{n=1}^{\infty}\frac{c'_{\Bar{\chi}}(n)}{n^{s}}.
\end{align}
\noindent
Let $\mathrm{Re}(s)=\sigma$, then the following inequality
\begin{align*}
|D_2(s,r)|
\leq\frac{\zeta(\sigma)\zeta(\sigma+r)\zeta(\sigma+r+1)}{\zeta(2(\sigma+r))}\zeta(r)E_r(\sigma)
\end{align*}
holds for $\sigma>1$, where
\begin{align}\label{erdefinition}
E_r(\sigma)
=C'(r)\prod_{p}(1+J_{p,r}(\sigma))
\end{align}
with
\begin{align*}
J_{p,r}(\sigma)
=-\frac{\frac{1-p^{-r}}{p^{r+1}}\left(\frac{1}{p^{\sigma}}+\frac{1}{2p^{2\sigma+2r+1}}\right)}{1+\frac{1-p^{-r}}{p^{r+1}}}
\end{align*}
and
\begin{align}\label{cprimerdefinition}
C'(r)
\coloneqq\prod_{p}\left(1+\frac{1}{p^{r+1}}\left(1-\frac{1}{p^r}\right)\right).
\end{align}
\end{lemma}

\begin{align}
D_2(s,r)
\coloneqq\sum_{m=1}^{\infty}\frac{1}{\varphi(m)m^{r+1}}\sum_{\substack{\chi(\bmod{m})\\\chi\ne\chi_0}}\tau(\chi)\sum_{n=1}^{\infty}\frac{c'_{\Bar{\chi}}(n)}{n^{s}}.
\end{align}

\begin{lemma}\label{shfiteddirichletseries}
For a fixed integer $r>1$, the Dirichlet series
\begin{align*}
A_{\sigma_r}(s)
=\sum_{n=1}^{\infty}\frac{\sigma_r(n+1)}{n^s}
=&\zeta(r+1)\sum_{i=0}^{r}\binom{r}{i}D_1(s-i,r)+\zeta(r+1)\sum_{i=0}^{r}\binom{r}{i}D_2(s-i,r),
%=&\zeta(r+1)\sum_{i=0}^{r}\frac{\zeta(s-i)}{\zeta(s+r+1-i)}K_r(s-i)+\zeta(r+1)\sum_{i=0}^{r}
\end{align*}
converges for $\mathrm{Re}(s)>1$, where $D_1(s,r)$ and $D_2(s,r)$ are defined in~\eqref{dskdefinition} and~\eqref{d2sksum}, respectively.
\end{lemma}

\begin{proof}
Employing the identity~\eqref{kthdivisoridentity}, we express
\begin{align*}
A_{\sigma_r}(s)
=\sum_{n=1}^{\infty}\frac{\sigma_r(n+1)}{n^s}
=\zeta(r+1)\sum_{n=1}^{\infty}\frac{(n+1)^{r}}{n^s}\sum_{m=1}^{\infty}\frac{c_m(n+1)}{m^{r+1}}.
\end{align*}
\noindent
Given that $r\geq 1$ is fixed, and under the conditions $\mathrm{Re}(s)>1$, and $\mathrm{Re}(s)+r>1$, one may interchange the order of summation, thereby obtaining
\begin{align}\label{sum1}
\zeta(r+1)\sum_{m=1}^{\infty}\frac{1}{m^{r+1}}\sum_{n=1}^{\infty}\frac{c_m(n+1)}{n^s}(n+1)^r
=\zeta(r+1)\sum_{i=0}^{r}\binom{r}{i}\sum_{m=1}^{\infty}\frac{1}{m^{r+1}}\sum_{n=1}^{\infty}\frac{c_m(n+1)}{n^{s-i}}.
\end{align}

Now applying Lemma~\ref{shiftedramanujansumformula} yields,
\begin{align*}
\sum_{m=1}^{\infty}\frac{1}{m^{r+1}}\sum_{n=1}^{\infty}\frac{c_m(n+1)}{n^{s-i}}
=&\sum_{m=1}^{\infty}\frac{1}{\varphi(m)m^{r+1}}\sum_{n=1}^{\infty}\frac{1}{n^{s-i}}\left(\mu(m)c_m(n)+\sum_{\substack{\chi(\bmod{m})\\ \chi\ne \chi_0}}c_{\chi}(1)c'_{\bar{\chi}}(n)\right)\\
=&\sum_{m=1}^{\infty}\frac{\mu(m)}{\varphi(m)m^{r+1}}\sum_{n=1}^{\infty}\frac{c_m(n)}{n^{s-i}}+\sum_{m=1}^{\infty}\frac{1}{\varphi(m)m^{r+1}}\sum_{\substack{\chi(\bmod{m})\\ \chi\ne \chi_0}}\tau(\chi)\sum_{n=1}^{\infty}\frac{c'_{\bar{\chi}}(n)}{n^{s-i}}\\
=&D_1(s-i,r)+D_2(s-i,r).
\end{align*}

\noindent
By applying Lemma~\ref{dsklemma} and Lemma~\ref{gausssumlemma}, in conjunction with~\eqref{sum1}, we thereby complete the proof.
\end{proof}

\end{document}
